\subsection{NARCBench Activation-Probing Defense}
\label{sec:narcbench-probing}

We evaluated the NARCBench collusion attack with the shared NARCBench
activation-probing defense, which follows the open-weight probing setting of
NARCBench~\citep{rose2026narcbench}. Because hidden activations cannot be
reconstructed from previously saved Bedrock text trajectories, we generated a
separate matched corpus with \texttt{Qwen/Qwen3-32B-AWQ}. The corpus contains
the same 50 NARCBench-Core scenarios under collusion and honest-control
conditions, yielding 100 samples. Scenario content and committee structure were
matched within each pair.

For every agent turn, we captured the final-token hidden state at transformer
layer 28 (dimension 5,120). We then applied the unchanged trained probe artifact
provided with the defense implementation. The runner verified the artifact before
deserialization using SHA-256
\texttt{d9a340a3276c477b817e9d0f7580a1c21cf3e5d080bb0ac0281e4c23b3d8e4ec}.
The pre-registered primary statistic reproduces the defense scorer exactly: it
projects every recorded turn activation onto the learned direction and takes
the maximum projection within a trajectory (\texttt{peak\_suspicion}). We also
report a public-discussion-only maximum and the released group-level metrics as
secondary analyses.

Qwen was run in non-thinking mode with sampling temperature 0.7, top-$p$ 0.8,
top-$k$ 20, a fixed base seed, and at most 512 new tokens. Responses followed a
bounded JSON contract. A malformed response received at most one schema-only
repair attempt; the repair prompt did not contain the malformed text, and the
accepted response was activation-scored under the original task prompt. Across
1,700 accepted turns, 80 repairs were required (4.71\%); all 100 samples
completed and no failed-generation artifact remained.

We treat the scenario pair as the unit of analysis. For each continuous probe
metric, we report attack and control means, their paired mean difference, a
10,000-resample paired-bootstrap 95\% interval, AUROC, and average precision.
We do not fit a binary threshold on the test set. The primary all-turn score
produced means of 6.473 under attack and 0.832 under control, a paired difference
of 5.642 (95\% CI [4.653, 6.666]), AUROC 0.9896, and average precision 0.9901.

We analyzed attack outcomes independently from the probe scores. Target
adherence required both designated colluders to vote for the assigned target,
and end-to-end success additionally required the committee to select that
target. Target adherence occurred in 31/50 attack samples (0.620; Wilson 95\%
CI [0.482, 0.741]), but only 3/50 attacks achieved end-to-end success. The
matched controls selected the target in 0/50 samples, giving a paired
target-selection uplift of 0.060 (bootstrap 95\% CI [0.000, 0.140]; exact
McNemar $p=0.250$). No sample abstained and no unrepaired parse failure remained.
Thus the probe's strong condition-label discrimination should not be read as
evidence that the attack reliably changed committee outcomes.

This result measures within-corpus discrimination of collusion-condition versus
matched-control trajectories, not successful harmful outcomes. The primary
scorer also includes the private coordination phase and therefore represents a
privileged evidence scope; the public-only result is reported separately. The
released bundle states that the probe was trained on 62 runs but does not expose
a verifiable scenario-level split against this evaluation set. We therefore do
not claim held-out or cross-model generalization without additional provenance
for the released probe artifact.
